% ============================================================================
% Fichier  : partie4/P04_C11_vocabulaire_risque.tex
% Titre    : Concepts et Vocabulaire du Risque
% Auteur   : MINKA MI NGUIDJOÏ Thierry Emmanuel
% ============================================================================
% Sources : ISO/IEC 27005:2022 ; ISO 31000:2018. Mairie d'Eséka : cas fictif.
% ============================================================================

\chapter{Concepts et Vocabulaire du Risque}
\label{chap:vocabrisque}
\chaptermark{Vocabulaire du risque}

\epigraph{%
    « On ne traite bien que les risques que l'on a su nommer. »%
}{Principe du cours}

\niveauChapitre{Fondamental}{La partie III.}

\begin{boxobjectifs}
\begin{itemize}
    \item Définir actif, menace, vulnérabilité, impact et probabilité.
    \item Énoncer les formules du risque et les interpréter.
    \item Distinguer risque inhérent et risque résiduel.
    \item Comprendre l'appétence et les critères d'acceptation du risque.
    \item Appliquer le vocabulaire à un cas de la mairie.
\end{itemize}
\end{boxobjectifs}

\bigskip

% ============================================================================
\section*{Question directrice : qu'est-ce, au juste, qu'un risque ?}
% ============================================================================

Le mot « risque » est employé à tort et à travers. Pour l'apprécier, il faut
le décomposer en éléments précis. Ce chapitre fixe le vocabulaire commun à
toutes les méthodes.

\medskip
\textbf{Plan du chapitre.}
\begin{enumerate}
    \item Les cinq notions (section~\ref{sec:c11-notions}).
    \item La formule du risque (section~\ref{sec:c11-formule}).
    \item Inhérent et résiduel (section~\ref{sec:c11-residuel}).
    \item Appétence et acceptation (section~\ref{sec:c11-appetence}).
    \item Un cas appliqué (section~\ref{sec:c11-cas}).
\end{enumerate}

\separateur

% ============================================================================
\section{Les Cinq Notions}
\label{sec:c11-notions}
% ============================================================================

\begin{boxdef}
\begin{itemize}
    \item \termecle{Actif}~: élément de valeur à protéger.
    \item \termecle{Menace}~: cause potentielle d'un incident indésirable.
    \item \termecle{Vulnérabilité}~: faiblesse exploitable par une menace.
    \item \termecle{Impact}~: conséquence d'un incident.
    \item \termecle{Probabilité} (ou vraisemblance)~: chance qu'un incident se
          produise.
\end{itemize}
\end{boxdef}

Ces notions s'enchaînent~: une \emph{menace} exploite une \emph{vulnérabilité}
d'un \emph{actif}, avec une certaine \emph{probabilité}, et produit un
\emph{impact}. Supprimer l'un des maillons, c'est réduire le risque.

% ============================================================================
\section{La Formule du Risque}
\label{sec:c11-formule}
% ============================================================================

\begin{boxdef}
Le \termecle{risque} se formule, selon le niveau de détail, par~:
\[
\text{Risque} = \text{Menace} \times \text{Vulnérabilité} \times \text{Impact}
\quad\text{ou}\quad
\text{Risque} = \text{Probabilité} \times \text{Impact}.
\]
\end{boxdef}

Ces formules ne sont pas des calculs exacts mais des \emph{grilles de
raisonnement}~: elles rappellent que le risque croît avec la probabilité et
avec l'impact, et qu'un impact nul (actif sans valeur) annule le risque.

% ============================================================================
\section{Inhérent et Résiduel}
\label{sec:c11-residuel}
% ============================================================================

\begin{boxdef}
\termecle{Risque inhérent}~: niveau de risque avant toute mesure.
\termecle{Risque résiduel}~: niveau qui subsiste après traitement. Le risque
résiduel n'est jamais nul~: il est accepté, explicitement, par la direction.
\end{boxdef}

Distinguer les deux évite deux illusions~: croire qu'on part de zéro (il y a
toujours un risque inhérent) et croire qu'on arrive à zéro (il reste toujours
un résiduel). Le management consiste à ramener le résiduel sous un seuil
accepté.

% ============================================================================
\section{Appétence et Acceptation}
\label{sec:c11-appetence}
% ============================================================================

\begin{boxdef}
\termecle{Appétence au risque}~: niveau de risque que l'organisation est prête
à accepter. Les \termecle{critères d'acceptation} traduisent cette appétence en
seuils concrets.
\end{boxdef}

Pour une mairie, l'appétence est faible sur l'état civil (qui fait foi) et plus
élevée sur le site web d'information. Fixer l'appétence est une décision de
gouvernance (le maire), pas un choix technique.

\begin{boxmetier}{RSSI}
\textbf{Chiffrer sans faux-semblant.} La RSSI sait que « Risque = P × I » ne
produit pas un nombre scientifique, mais un classement défendable. L'important
n'est pas la précision du chiffre~: c'est la cohérence du classement et la
traçabilité du raisonnement.
\end{boxmetier}

% ============================================================================
\section{Un Cas Appliqué}
\label{sec:c11-cas}
% ============================================================================

\begin{boxobjet}{Le registre d'état civil}
\textbf{Constatation.} Actif~: registre d'état civil numérisé. Menace~:
rançongiciel apporté par clé \texttt{USB}. Vulnérabilité~: pas de sauvegarde,
accès partagé. Impact~: service à l'arrêt, actes indisponibles.

\medskip
\textbf{Analyse} (confiance élevée). Probabilité élevée (la menace s'est déjà
réalisée), impact majeur~: risque inhérent très élevé. Deux mesures,
la sauvegarde quotidienne et l'interdiction des clés personnelles, réduisent la
probabilité et l'impact, abaissant le risque résiduel à un niveau acceptable.
\end{boxobjet}

\separateur

\begin{boxresume}
\textbf{Ce qu'il faut retenir de ce chapitre~:}
\begin{itemize}
    \item Un risque enchaîne \textbf{actif, menace, vulnérabilité,
          probabilité, impact}.
    \item Les \textbf{formules} du risque sont des grilles de raisonnement, non
          des calculs exacts.
    \item Le \textbf{risque résiduel} subsiste toujours et s'accepte
          explicitement.
    \item L'\textbf{appétence} au risque est une décision de gouvernance.
    \item Le vocabulaire s'applique à chaque \textbf{actif} de la mairie.
\end{itemize}
\end{boxresume}

\begin{boxinfo}
\textbf{Points de contrôle de l'élève.} Je sais~: définir les cinq notions~;
énoncer les formules du risque~; distinguer inhérent et résiduel~; expliquer
l'appétence~; analyser un cas.
\end{boxinfo}

\bigskip

\begin{boxexo}[title={\ding{45}\quad Exercice~11.1. Les cinq notions \niveaubase}]
Définissez actif, menace, vulnérabilité, impact et probabilité.
\end{boxexo}

\begin{boxexo}[title={\ding{45}\quad Exercice~11.2. La chaîne \niveaubase}]
Pour le vol du fichier des contribuables, nommez la menace, la vulnérabilité et
l'impact.
\end{boxexo}

\begin{boxexo}[title={\ding{45}\quad Exercice~11.3. Inhérent ou résiduel \niveaubase}]
Un risque « avant toute mesure » est-il inhérent ou résiduel~?
\end{boxexo}

\begin{boxexo}[title={\ding{45}\quad Exercice~11.4. Formule \niveaumoyen}]
Pourquoi un impact nul annule-t-il le risque, même si la probabilité est
élevée~?
\end{boxexo}

\begin{boxexo}[title={\ding{45}\quad Exercice~11.5. Résiduel jamais nul \niveaumoyen}]
Expliquez pourquoi le risque résiduel n'est jamais nul.
\end{boxexo}

\begin{boxexo}[title={\ding{45}\quad Exercice~11.6. Appétence \niveaumoyen}]
Pourquoi l'appétence au risque de la mairie diffère-t-elle entre l'état civil
et le site web~?
\end{boxexo}

\begin{boxexo}[title={\ding{45}\quad Exercice~11.7. Casser un maillon \niveaumoyen}]
Montrez qu'agir sur la vulnérabilité réduit le risque, sur l'exemple de la
mairie.
\end{boxexo}

\begin{boxexo}[title={\ding{45}\quad Exercice~11.8. Chiffre défendable \niveaumoyen}]
« Risque = P × I donne un nombre scientifique. » Nuancez.
\end{boxexo}

\begin{boxexo}[title={\ding{45}\quad Exercice~11.9. Analyse d'actif \niveauexpert}]
Analysez le risque sur la messagerie de la mairie (actif, menace,
vulnérabilité, probabilité, impact, inhérent, résiduel).
\end{boxexo}

\begin{boxexo}[title={\ding{45}\quad Exercice~11.10. Critères d'acceptation \niveauexpert}]
Proposez des critères d'acceptation du risque pour la mairie d'Eséka.
\end{boxexo}

\bigskip

\begin{boxinfo}
\textbf{Pour aller plus loin.}
\begin{itemize}
    \item ISO/IEC 27005~\cite{iso27005}~; ISO 31000~\cite{iso31000}.
\end{itemize}
\end{boxinfo}

\begin{boxinfo}
\textbf{Transition.} Le vocabulaire posé, le Chapitre~\ref{chap:methodes}
présente les méthodes qui l'appliquent~: EBIOS RM, OCTAVE, MEHARI.
\end{boxinfo}
